% Problem record op_68d6ff4ef7be1577. This TeX file is derived from
% database/problems_json/op_68d6ff4ef7be1577.json by scripts/sync-tex.mjs; edit the JSON record
% and rerun the script rather than editing this file.
\section{Two-way distillable entanglement of mixed two-mode Gaussian states}
\label{sec:op_68d6ff4ef7be1577}

\paragraph{Problem.}
What is the exact two-way distillable entanglement of a finite-energy mixed two-mode Gaussian state?

Let $\rho_{AB}$ be a finite-energy mixed two-mode Gaussian density operator. Let $\Lambda_k$ range over protocols using local operations and unlimited two-way classical communication. The complete protocols are trace preserving, and non-Gaussian operations are allowed. Finite energy means finite total mean photon number. Alice and Bob hold one mode each. Define the Bell-pair density operator and trace norm by

\begin{equation}
\Phi_2:=\tfrac12(|00\rangle+|11\rangle)(\langle00|+\langle11|),
\qquad
\lVert X\rVert_1:=\operatorname{Tr}\sqrt{X^\dagger X}.
\label{eq:ser13-statement-1}
\end{equation}

Using the Bell pair in Eq.~\eqref{eq:ser13-statement-1}, call $r\geq0$ achievable if an allowed sequence $(\Lambda_k)_{k\geq1}$ satisfies

\begin{equation}
\lim_{k\to\infty}\left\lVert\Lambda_k(\rho_{AB}^{\otimes k})-\Phi_2^{\otimes\lfloor rk\rfloor}\right\rVert_1=0.
\label{eq:ser13-statement-2}
\end{equation}

Determine $D_{\leftrightarrow}(\rho_{AB}):=\sup\{r\geq0:r\text{ is achievable}\}$, with achievability defined by Eq.~\eqref{eq:ser13-statement-2}.

\subsection*{Status}

Unsolved

\subsection*{Source}

This exact-rate question is formulated from the gap between quantitative distillation bounds \sourcecite{ref:ser13-1}{DW05}, \sourcecite{ref:ser13-2}{VW02}, together with the Gaussian distillability criterion \sourcecite{ref:ser13-3}{GDCZ01}.

\subsection*{Progress}

\begin{itemize}
  \item Giedke et al., Theorem 1, prove that a bipartite Gaussian state is distillable exactly when its partial transpose is not positive. In particular, every entangled two-mode Gaussian state is distillable. This settles positivity of the rate, not its exact value. \sourcecite{ref:ser13-3}{GDCZ01}

  \item Devetak--Winter, Theorem 10, proves the finite-dimensional hashing inequality. The hashing bound, applied in either communication direction and extended to finite-energy bosonic states by local truncation, gives

\begin{equation}
D_{\leftrightarrow}(\rho_{AB})\geq\max\{0,S(\rho_A)-S(\rho_{AB}),S(\rho_B)-S(\rho_{AB})\},
\label{eq:ser13-progress-2-1}
\end{equation}

In Eq.~\eqref{eq:ser13-progress-2-1}, $\rho_A:=\operatorname{Tr}_B\rho_{AB}$, $\rho_B:=\operatorname{Tr}_A\rho_{AB}$, and $S(\omega):=-\operatorname{Tr}(\omega\log_2\omega)$. \sourcecite{ref:ser13-1}{DW05}

  \item Vidal--Werner, Proposition 7 and Section V.C, bounds the rate by logarithmic negativity:

\begin{equation}
D_{\leftrightarrow}(\rho_{AB})\leq\log_2\lVert\rho_{AB}^{T_B}\rVert_1
=\sum_{j=1}^2\max\{0,-\log_2\widetilde\nu_j\},
\label{eq:ser13-progress-3-1}
\end{equation}

In Eq.~\eqref{eq:ser13-progress-3-1}, $T_B$ is partial transposition and $\widetilde\nu_j$ are the symplectic eigenvalues of the partially transposed covariance in vacuum-$I$ units. \sourcecite{ref:ser13-2}{VW02}
\end{itemize}

\subsection*{References}

\begin{enumerate}[label={},leftmargin=4.8em,labelsep=0.5em]
  \footnotesize
  \item[\textup{[DW05]}]\label{ref:ser13-1}
  I. Devetak and A. Winter, "Distillation of Secret Key and Entanglement from Quantum States," \emph{Proceedings of the Royal Society A: Mathematical, Physical and Engineering Sciences} \textbf{461}, 207--235 (2005). \href{https://doi.org/10.1098/rspa.2004.1372}{doi:10.1098/rspa.2004.1372}; \href{https://arxiv.org/abs/quant-ph/0306078}{arXiv:quant-ph/0306078}.

  \item[\textup{[VW02]}]\label{ref:ser13-2}
  G. Vidal and R. F. Werner, "Computable Measure of Entanglement," \emph{Physical Review A} \textbf{65}, 032314 (2002). \href{https://doi.org/10.1103/PhysRevA.65.032314}{doi:10.1103/PhysRevA.65.032314}; \href{https://arxiv.org/abs/quant-ph/0102117}{arXiv:quant-ph/0102117}.

  \item[\textup{[GDCZ01]}]\label{ref:ser13-3}
  G. Giedke, L.-M. Duan, J. I. Cirac, and P. Zoller, "Distillability Criterion for All Bipartite Gaussian States," \emph{Quantum Information and Computation} \textbf{1}(3), 79--86 (2001). \href{https://arxiv.org/abs/quant-ph/0104072}{arXiv:quant-ph/0104072}.
\end{enumerate}

\subsection*{Comment}

Separable states have rate zero. For general entangled mixed two-mode Gaussian states, the known lower and upper bounds need not coincide. Gaussian-operation no-go theorems do not settle the unrestricted LOCC rate.

\subsection*{Field}

Quantum Resource Theory

\subsection*{Topic}

Gaussian quantum information; Entanglement distillation

\subsection*{ID}

\texttt{op\_68d6ff4ef7be1577}
